\documentclass[12pt]{article}
\usepackage{amsmath,amssymb,amsfonts}
\usepackage[margin=1in]{geometry}

\begin{document}

\begin{center}
{\Large\bfseries Chapter 10, Problem 18}
\end{center}

\noindent\textbf{Problem.} Using the results of Problem 10.4, find the characters and representations of the group of rotational symmetries of the regular tetrahedron.

\bigskip
\noindent\textbf{Solution.}

From Problem 4, the rotation group of the regular tetrahedron is $A_4$ with 12 elements in 4 conjugacy classes:
\begin{align*}
\mathcal{C}_1 &= \{e\} &&(1 \text{ element}), \\
\mathcal{C}_2 &= \{(123), (142), (134), (243)\} &&(4 \text{ elements}), \\
\mathcal{C}_3 &= \{(132), (143), (124), (234)\} &&(4 \text{ elements}), \\
\mathcal{C}_4 &= \{(12)(34), (13)(24), (14)(23)\} &&(3 \text{ elements}).
\end{align*}

With 4 classes, there are 4 irreducible representations. Since $\sum n_i^2 = 12$ and $n_i \geq 1$:
\[
1^2 + 1^2 + 1^2 + 3^2 = 12.
\]
So there are three 1-dimensional and one 3-dimensional irreducible representations.

\medskip
\noindent\textbf{One-dimensional representations:}

For 1D representations, $D$ is a homomorphism from $A_4$ to $\mathbb{C}^*$. Since elements of $\mathcal{C}_4$ have order 2, $\chi(\mathcal{C}_4) = \pm 1$. Since elements of $\mathcal{C}_2, \mathcal{C}_3$ have order 3, $\chi(\mathcal{C}_2)^3 = 1$ and $\chi(\mathcal{C}_3)^3 = 1$.

Let $\omega = e^{2\pi i/3}$. The three 1D representations are:

\[
\begin{array}{c|cccc}
 & \mathcal{C}_1 & \mathcal{C}_2 & \mathcal{C}_3 & \mathcal{C}_4 \\
 & 1 & 4 & 4 & 3 \\
\hline
\Gamma_1 & 1 & 1 & 1 & 1 \\
\Gamma_2 & 1 & \omega & \omega^2 & 1 \\
\Gamma_3 & 1 & \omega^2 & \omega & 1
\end{array}
\]

Note: $\chi(\mathcal{C}_4) = 1$ for all 1D representations because $(12)(34) = (12)(34)$ can be written as a product of elements whose images must give 1 (since $(12)(34) \cdot (13)(24) = (14)(23)$ and the constraint from the group structure forces this).

\medskip
\noindent\textbf{Three-dimensional representation:}

The character $\chi^{(4)}$ can be found from the orthogonality relations. Let $\chi^{(4)} = (3, a, b, c)$.

Orthogonality with $\Gamma_1$:
\[
1\cdot 3 + 4\cdot a + 4\cdot b + 3\cdot c = 0 \implies 4a + 4b + 3c = -3.
\]

Orthogonality with $\Gamma_2$:
\[
3 + 4\omega a + 4\omega^2 b + 3c = 0.
\]

Orthogonality with $\Gamma_3$:
\[
3 + 4\omega^2 a + 4\omega b + 3c = 0.
\]

Adding the last two equations: $6 + 4(\omega + \omega^2)(a + b) + 6c = 0$. Since $\omega + \omega^2 = -1$:
\[
6 - 4(a + b) + 6c = 0 \implies 4(a+b) - 6c = 6.
\]

From the first equation: $4(a+b) + 3c = -3$. Subtracting: $-9c = 9$, so $c = -1$.

Then $4(a+b) = -3 + 3 = 0$, so $a + b = 0$, i.e., $b = -a$.

Subtracting the $\Gamma_2$ and $\Gamma_3$ orthogonality equations:
\[
4(\omega - \omega^2)a + 4(\omega^2 - \omega)b = 0 \implies 4(\omega - \omega^2)(a - b) = 0.
\]
Since $\omega \neq \omega^2$: $a = b$. Combined with $b = -a$: $a = b = 0$.

\medskip
\noindent\textbf{Complete character table:}

\[
\begin{array}{c|cccc}
 & \{e\} & \{(123),\ldots\} & \{(132),\ldots\} & \{(12)(34),\ldots\} \\
 & 1 & 4 & 4 & 3 \\
\hline
\Gamma_1 & 1 & 1 & 1 & 1 \\
\Gamma_2 & 1 & \omega & \omega^2 & 1 \\
\Gamma_3 & 1 & \omega^2 & \omega & 1 \\
\Gamma_4 & 3 & 0 & 0 & -1
\end{array}
\]

where $\omega = e^{2\pi i/3}$.

\medskip
\noindent\textbf{Explicit 3D representation:}

The 3D representation can be constructed from the permutation representation. Assign vertex $k$ to basis vector $\mathbf{e}_k$ ($k = 1,2,3,4$). The permutation representation is 4-dimensional. It decomposes as $\Gamma_1 \oplus \Gamma_4$: the invariant subspace spanned by $\mathbf{e}_1 + \mathbf{e}_2 + \mathbf{e}_3 + \mathbf{e}_4$ gives $\Gamma_1$, and the orthogonal complement (the 3D subspace with $\sum x_i = 0$) gives $\Gamma_4$.

\end{document}
